% ===== tr-06 块1：第 4 章 pp.61--64（原文页 61--64）=====
% [原文 p.61]
\begin{example}{4.11.5}
回到例 4.11.3，对 $\ell=1,2,3$ 的情形，考虑附着于斜率 1 的射线的函数 $f_\mathfrak{d}$。在每种情形下，曲面 $X$ 都是 $\PP^2$，其坐标轴为 $D_1$、$D_2$ 与 $D_{\mathrm{out}}$。于是 $\tilde X$ 由在 $D_1$ 与 $D_2$ 上各爆破 $\ell$ 个点得到。

先考虑 $\ell=1$ 的情形。注意对 $\beta=dH$（即 $\PP^2$ 中 $d$ 次曲线的类），$\mathbf{P}$ 的唯一相关取法是 $\mathbf{P}_1=d$、$\mathbf{P}_2=d$，于是我们有
\[
\beta_{\mathbf{P}}=d\nu^*H-dE_{11}-dE_{21}.
\]
这代表的是一条 $d$ 次曲线、恰好 $d$ 次经过这两个爆破点的类。容易看出，这样一条曲线的唯一选择是过这两点的一条直线的 $d$ 重覆盖。此外，该覆盖必须在 $D_{\mathrm{out}}$ 上完全分歧，以保证与 $D_{\mathrm{out}}$ 的相切阶符合要求。这需要一个虚拟计数，相关的局部化计算在 [28] 中完成，给出 $N_{\mathbf{P}}=(-1)^{d+1}/d^2$。于是我们得到
\[
\log f_\mathfrak{d}=\sum_{d=1}^{\infty}d\left(\frac{(-1)^{d+1}}{d^{2}}\right)t^{2d}x^{-d}y^{-d}.
\]
取指数，可得 $f_\mathfrak{d}=1+t^{2}x^{-1}y^{-1}$，与例 4.11.3 一致。所以这里我们只是在数 $\PP^2$ 中过两点的唯一一条直线以及该直线的某些多重覆盖。
% [Tutor Note: $\mathbf{P}_1=d$、$\mathbf{P}_2=d$ 是“长度 1 的分拆 $(d)$”的简写；代入定义 $\beta_{\mathbf P}=\nu^*\beta-\sum_i\sum_j p_{ij}E_{ij}$ 时，$i=1,2$ 各只有 $j=1$ 一项、$p_{i1}=d$，故得 $d\nu^*H-dE_{11}-dE_{21}$。]

转到 $\ell=2$、$\beta=dH$ 的情形。当 $d=1$ 时，分拆有四种取法：$\mathbf{P}=(1+0,1+0)$、$(1+0,0+1)$、$(0+1,1+0)$ 与 $(0+1,0+1)$。每一种对应在 $D_1$、$D_2$ 上各选一个点，而通过每一对这样的点都有唯一一条直线。于是对每种这样的 $\mathbf{P}$ 都有 $N_{\mathbf{P}}=1$。与 $\ell=1$ 的情形一样，这些直线也通过多重覆盖对更高次数作出贡献：比如说，$\mathbf{P}=(d+0,d+0)$ 贡献 $N_{\mathbf{P}}=(-1)^{d+1}/d^{2}$。当 $d=2$ 时，可以看到对 $\mathbf{P}=(2+0,1+1)$ 之类没有曲线，因为这要求一条在 $D_1$ 上有结点、且与 $D_{\mathrm{out}}$ 相切的二次曲线；这样的曲线不存在。但对于 $\mathbf{P}=(1+1,1+1)$，我们看的是经过全部四点并与 $D_{\mathrm{out}}$ 相切的二次曲线。很容易看出这样的二次曲线有两条。

然后可以验证，唯一其他有贡献的曲线是四条直线之一或两条二次曲线的多重覆盖。二次曲线的多重覆盖贡献其实与直线不同，因为与 $D_{\mathrm{out}}$ 的相切阶不同。结果，二次曲线的 $d$ 重覆盖的正确贡献是 $1/d^{2}$。于是可得
\[
\log f_\mathfrak{d}=4\sum_{d=1}^{\infty}d\left(\frac{(-1)^{d+1}}{d^{2}}\right)t^{2d}x^{-d}y^{-d}+2\sum_{d=1}^{\infty}2d\left(\frac{1}{d^{2}}\right)t^{4d}x^{-2d}y^{-2d},
\]
取指数便得到
\[
f_\mathfrak{d}=\frac{(1+t^{2}xy)^{4}}{(1-t^{4}x^{-2}y^{-2})^{4}}=(1-t^{2}x^{-1}y^{-1})^{-4}.
\]
% [Tutor Note: 此处疑似原文笔误。分子按上下文应为 $(1+t^{2}x^{-1}y^{-1})^{4}$，原书排作 $(1+t^{2}xy)^{4}$（页 61，x 60--570pt、y 700--745pt 区域 900dpi 裁图 c61_a.png 可辨：“xy”上无指数）。判据：由上一式取指数，$\exp\bigl(4\sum_{d\ge1}(-1)^{d+1}u^{d}/d\bigr)=(1+u)^{4}$，$u=t^{2}x^{-1}y^{-1}$；又 $(1+u)^{4}/(1-u^{2})^{4}=(1-u)^{-4}$，与式中最后一个等号完全吻合。正文照录原样，此注给出正确形式。]
\end{example}
% [原文 p.62]
在 $\ell=3$ 的情形，人们预期有 $3\times3=9$ 条直线，因为对 $D_1$ 上一个点与 $D_2$ 上一个点的每一种选法，都有一条经过它们的直线。对于二次曲线，有这些直线的二重覆盖，贡献为 $-9/4$，以及 $2\times3\times3=18$ 条二次曲线。这里需要在 $D_1$ 上选两个点、在 $D_2$ 上选两个点，然后有两条经过这四个点且与 $D_{\mathrm{out}}$ 相切的二次曲线。

对于三次曲线，有直线的三重覆盖的贡献，总计 $9/9$，还有若干来自平面三次曲线的贡献。结果表明，对 $\mathbf{P}=(1+1+1,1+1+1)$，$N_{\mathbf{P}}=18$。注意这给出了经过 6 个固定点、且使 $D_{\mathrm{out}}$ 为其三重切线的结点平面三次曲线的计数。另一方面，对 $\mathbf{P}=(1+2+0,1+1+1)$，$N_{\mathbf{P}}=3$。注意这种形状的分拆总共有 12 个。后一个计数表示结点位于所选点之一、同时还经过另外四个所选点、且使 $D_{\mathrm{out}}$ 为三重切线的结点三次曲线。于是可以断言
\[
\log f_\mathfrak{d}=9t^{2}x^{-1}y^{-1}+2(-9/4+18)t^{4}x^{-2}y^{-2}+3(9/9+54)t^{6}x^{-3}y^{-3}+\cdots.
\]
与例 4.11.3 中给出的值直接比较，二者一致。

我们以对定理 4.11.4 的简短补充动机以及关于其证明的几句话来结束本节。

设我们有一块整仿射流形，如 图 4.11 所示。这里我们设想一个带有两个奇点的曲面，绕奇点的局部单值（在水平与竖直线段上）在适当选取的基（对每个线段不同）下分别为
\[
\begin{pmatrix}1&\ell_1\\0&1\end{pmatrix}
\quad\text{与}\quad
\begin{pmatrix}1&\ell_2\\0&1\end{pmatrix}.
\]
这比第 4.10 节所考虑的情形略为一般，那里我们只讨论单值形如
\[
\begin{pmatrix}1&1\\0&1\end{pmatrix}
\]
的奇点。尽管如此，那一节的技巧仍然适用，只是附着于从奇点射向中心顶点 $v$ 的初始射线的函数可以取为 $(1+x^{-1})^{\ell_1}$ 与 $(1+y^{-1})^{\ell_2}$ 的形式。这大致就是上面所讨论例子的形状。应用散射程序便会给出 $\check X_0(B,\mathscr{P},s)$ 的一个光滑化。然而，在镜像一侧，我们把 $B$ 解释为对偶相交复形，这意味着它应当来自一个退化 $\mathcal{X}\to D$，其中心纤维 $\mathcal{X}_0$ 有一个同构于 $\PP^2$ 的不可约分支 $Y_v$（对应顶点 $v$）。此外，全空间 $\mathcal{X}$ 应当有 $\ell_1+\ell_2$ 个位于 $Y_v$ 的环面边界上的通常二重点。若在 $\mathcal{X}$ 内部爆破 Weil 除子 $Y_v$，就得到这些通常二重点的一个小消解 $\tilde{\mathcal{X}}\to\mathcal{X}$，特别地，$Y_v$ 的真变换 $\tilde Y_v$ 就是 $Y_v$ 在点 $Y_v\cap\operatorname{Sing}(\mathcal{X})$ 处的爆破。这一操作在 $Y_v$ 的一条坐标轴上爆破 $\ell_1$ 个点、在另一条上爆破 $\ell_2$ 个点。这恰好就是定理 4.11.4 中所考虑的曲面。
% [Tutor Note: 这里的检查帽记号是 $\check X_0(B,\mathscr{P},s)$（$\check{}$ 加在 $X$ 上，第二个参数是 rsfs 花体 $\mathscr{P}$，页 62 字体为 rsfs10）；而 p.62 末段起的退化族与中心纤维分别用 $\mathcal{X}$、$\mathcal{X}_0$、$\tilde{\mathcal{X}}$（CMSY 花体加波浪号），两者不可混用。]
% [原文 p.63]
\figimg[0.72]{p4_fig11.png}
\figtitle{4.11}{}
% [Tutor Note: 原文图注只有“Figure 11.”，无描述性文字；图内只有两个 × 记号（奇点）与中心顶点标签 $v$（页图经 1400dpi 复核为普通斜体 $v$，其上方没有波浪号），对应正文“从奇点射向中心顶点 $v$ 的初始射线”。译文图注留空以忠实原文。]

下面考虑定理 4.11.4 所计数的 $\tilde Y_v$ 上那类曲线。它们是 $\tilde Y_v$ 中只在一个点与第三条坐标轴相交的曲线。可以把它们看作 $\tilde{\mathcal{X}}_0$ 中的曲线，但不是那种能形变到族 $\tilde{\mathcal{X}}\to D$ 的一般纤维中的全纯曲线的曲线。相反，粗略地说，我们期望这样的曲线形变为全纯圆盘：其中与 $\tilde{\mathcal{X}}_0$ 的奇异轨迹的交点（即与 $\tilde Y_v$ 的第三轴的交点）膨胀成一个 $S^1$，给出全纯圆盘的边界。大致地，我们期望这个边界落在族 $\tilde{\mathcal{X}}\to D$ 的一般纤维上某个 SYZ 纤维化的纤维中。该边界在纤维中的同调类由曲线与第三轴的相切阶决定。

定理 4.11.4 中所考虑的相对曲线（与上述全纯圆盘）之间的这一对应只在“道义”的层面上成立；目前还没有证明我们真的在数这样的全纯圆盘。然而，这一论证是定理 4.11.4 的主要动机。

最后，就证明而言，有几个步骤。首先，我们证明散射图可以形变得看起来像热带曲线的并，并用 Mikhalkin 的基本曲线计数结果 [62]（由 Nishinou 与 Siebert [64] 发展的版本）来证明散射图在环面曲面上实现某些曲线计数。然后，用 Jun Li 的粘合公式 [57] 把这与爆破曲面的 Gromov--Witten 计数联系起来。

\section{其他近期结果与未来}

我将以对这里所讨论方法的应用与未来发展的简短讨论来收尾。
% [原文 p.64]
最近，这里所描述的光滑化机制的一个变体被我自己、Hacking 与 Keel [26] 用来给出对 $(Y,D)$ 这类对的镜像的一个非常一般的构造，其中 $Y$ 是有理曲面、$D$ 是构成有理曲线环的有效反典范除子。我们利用 [28] 写下我们所称的\emph{典范散射图}，它可以完全用对 $(Y,D)$ 的 Gromov--Witten 理论（更具体地说，是与 $D$ 只交于一个点的曲线计数）来描述。这一散射图决定了镜像族。然而，还用到了一件额外的关键工具，以把所构造的族部分紧化。这是必要的，因为与本文所考虑的仿射流形不同，与 $(Y,D)$ 自然相伴的那个仿射流形在多面体分解的一个顶点处有奇点。在这个顶点处没有光滑化的局部模型，因此在构造族时，所得到的族会“缺少”一个点。要把这个点加回来，就需要确保所构造的族上有足够多的函数，而事实证明同调镜像对称为构造这样的函数提示了一种自然的办法。这可以热带地完成，从而造出我们所称的 theta 函数。把同样的构造应用于退化阿贝尔簇的情形，确实产生通常的 theta 函数；我们预期在别的这些场合所构造的函数也会同样有用。关于这些想法的一个综述见 [35]。
% [Tutor Note: “典范散射图”原文为 canonical scattering diagram；它与本文 §11 由 $\mathfrak D$ 出发经散射程序得到的 $S(\mathfrak D)$ 不是同一个东西------后者是对给定散射图的完备化，而前者在这里是由 $(Y,D)$ 的 Gromov--Witten 理论直接写出的对象，用来确定镜像族。术语按 glossary 的“散射图”处理。]

[26] 的构造还解决了一个早于镜像对称的问题。特别地，我们证明了 Looijenga 关于尖点奇点可光滑性的一个猜想。

可以把 theta 函数看作仿射簇上函数环的典范基，或射影簇上线丛截影空间的典范基。因此，它们与 cluster 代数理论中的典范基发生联系，为构造 cluster 代数的典范基提供了一个框架。

我们还预期，对曲面组 $(Y,D)$ 的技巧将会推广。事实上，K3 曲面的任何极大幂零正规交叉退化的镜像伴侣都可以沿类似思路构造，这是与 Hacking、Keel 及 Siebert 的合作中的工作。预期是：再加上对数 Gromov--Witten 理论的帮助，人们应当能够对 Calabi--Yau 流形的极大幂零退化的镜像伴侣，在所有维数上写出一个一般的构造。

仍然遗留的问题是：从周期中提取计数信息------正是这一点最初给了镜像对称以激动人心之处。这里我们展示了计数几何如何可以反映在镜像中，但只是以一种相当局部的方式。我们预期应当能够进行周期积分的计算，以提取镜像的亏格零 Gromov--Witten 不变量，但在这个方向上还有一些技术问题。尽管如此，通过退化来理解镜像对称的这一纲领（它受 SYZ 猜想启发）似乎提供了一个强有力的框架，使人得以在代数几何与热带几何的领域内思考镜像对称。
% [Tutor Note: 本段“周期”与“周期积分”原文为 period、period integrals，指全纯体积形式的周期；末句“亏格零 Gromov--Witten 不变量”即 genus zero 情形，与 §11 中只涉及有理（亏格零）曲线一致。另注意原文“one should be able write down”漏了 to（p.64，页图可辨），属原文排印小疵，译文按正常语法处理。]
